% !TEX program = xelatex
% ============================================================
% SLB 技术文档：6.5.6 调制方式
% 规范：slb-doc-generator / tech-doc-style-guide.md
% 模块类型：通用功能
% ============================================================
\documentclass[11pt,a4paper]{ctexart}

\usepackage{amsmath,amssymb,amsfonts}
\usepackage{booktabs}
\usepackage{geometry}
\usepackage{hyperref}
\usepackage{enumitem}
\usepackage{xcolor}
\usepackage{longtable}
\usepackage{array}

\geometry{left=2.5cm,right=2.5cm,top=2.5cm,bottom=2.5cm}
\hypersetup{colorlinks=true,linkcolor=blue,urlcolor=blue}

\newcolumntype{L}[1]{>{\raggedright\arraybackslash}p{#1}}

\title{%
  \textbf{星闪 SLB 物理层}\\[0.4em]
  \Large 6.5.6 调制方式技术文档\\[0.3em]
  \large QPSK \textbullet\ 16QAM\textasciitilde 4096QAM \textbullet\ Gray 嵌套映射
}
\author{依据 T/XS 10001-2025《星闪无线通信系统 接入层\\
同步低时延宽带空口 SLB 技术要求和测试方法》整理\\
（团体标准 20250326 版）}
\date{2026 年 7 月 27 日}

\begin{document}
\maketitle
\tableofcontents
\newpage

%======================================================================
\part{总览}
%======================================================================

\section{文档范围与模块划分}

本文档对应 T/XS 10001-2025 第 6.5.6 节``调制方式''，规定编码/加扰后比特序列 $b(k)$ 到复数符号 $d(i)$ 的映射。只负责\textbf{比特域$\rightarrow$符号域}；CRC、加扰、RE 映射与功控不在本节。

\begin{longtable}{L{2.2cm}L{3.2cm}L{2cm}L{6.5cm}}
\caption{文档范围与模块划分\label{tab:656-scope}} \\
\toprule
\textbf{子节号} & \textbf{调制} & \textbf{$m$} & \textbf{归一化} \\
\midrule
6.5.6.1 & QPSK & 2 & $1/\sqrt{2}$ \\
6.5.6.2 & 16QAM & 4 & $1/\sqrt{10}$ \\
6.5.6.3 & 64QAM & 6 & $1/\sqrt{42}$ \\
6.5.6.4 & 256QAM & 8 & $1/\sqrt{170}$ \\
6.5.6.5 & 1024QAM & 10 & $1/\sqrt{682}$ \\
6.5.6.6 & 4096QAM & 12 & $1/\sqrt{2730}$ \\
\bottomrule
\end{longtable}

\section{与相关章节的关系}

\begin{longtable}{L{2.2cm}L{1.8cm}L{9.5cm}}
\caption{与相关章节的关系\label{tab:656-rel}} \\
\toprule
\textbf{相关节号} & \textbf{关系} & \textbf{说明} \\
\midrule
6.2.5 & 上游 & MCS 表决定调制阶；第一/二类数据 TB 链路 \\
6.2.4 & 上游 & 多数控制固定 QPSK；TCI 二阶段 ACK 可为 QPSK/16QAM \\
6.2.6 & 上游 & 信道编码输出比特 \\
6.5.7 & 上游 & 数据面必须先加扰再调制；输入为 $\tilde{g}_k$ \\
6.5.3 & 平行 & 参考信号 QPSK（无前置负号），\textbf{禁止}与 6.5.6.1 混用接口 \\
6.2.3 & 下游 & 资源映射消费 $d(i)$ \\
6.5.4 & 下游 & 平均能量为 1 的符号再按功控缩放 \\
6.5.8 & 下游 & 深覆盖路径：调制后可进 SC-FDMA \\
\bottomrule
\end{longtable}

\section{通用功能端到端流水线}

\begin{center}
\small
\textbf{控制}：编码 $\rightarrow$ \textbf{6.5.6.1 QPSK} $\rightarrow$ RE 映射\\[0.3em]
\textbf{数据}：CRC/编码 $\rightarrow$ \textbf{6.5.7} 加扰 $\rightarrow$ \textbf{6.5.6.x} $\rightarrow$（可选 6.5.8）$\rightarrow$ 映射 $\rightarrow$ 6.5.4 功控\\[0.3em]
\textbf{禁混}：参考信号走 6.5.3，不走本节调制 API
\end{center}

\newpage
%======================================================================
\part{6.5.6 调制方式（工程解读）}
%======================================================================

\section{协议章节对应}

本节对应 T/XS 10001-2025 第 6 章第 6.5.6 节，涵盖 6.5.6.1\textasciitilde 6.5.6.6 六种调制。

\section{功能定义}

\begin{enumerate}[nosep]
  \item \textbf{映射}：按 MCS/配置选定的调制阶，将 $m$ 比特映射为 1 个复数符号 $d(i)$；
  \item \textbf{归一化}：星座平均能量为 1（$1/\sqrt{N_{\mathrm{avg}}}$）；
  \item \textbf{Gray 嵌套}：16QAM 及以上按 $2-(1-2b)$ 嵌套生成 I/Q 幅度；
  \item \textbf{强制前置负号}：标准公式均含整体负号，实现不得省略。
\end{enumerate}

\section{模块边界（上下游及职责划分）}

\subsection{上游模块}

\begin{longtable}{L{4.2cm}L{4.5cm}L{5.5cm}}
\caption{上游模块输入\label{tab:656-up}} \\
\toprule
\textbf{上游} & \textbf{输入} & \textbf{说明} \\
\midrule
信道编码（6.2.6） & 控制比特 $g_k$ & 多数控制直接 QPSK \\
比特加扰（6.5.7） & 数据 $\tilde{g}_k$ & 数据面必经加扰 \\
MCS/调度 & \texttt{mod\_type} & 表 8\textasciitilde 10 或固定 QPSK \\
\bottomrule
\end{longtable}

\subsection{下游模块}

\begin{longtable}{L{4.2cm}L{4.5cm}L{5.5cm}}
\caption{下游模块输出\label{tab:656-down}} \\
\toprule
\textbf{下游} & \textbf{输出} & \textbf{说明} \\
\midrule
资源映射（6.2.3） & $d(i)$ & 写入 RE \\
SC-FDMA（6.5.8） & $d(i)$ 作 $x(n)$ & 仅 6.3.3 深覆盖 \\
功控（6.5.4） & 单位能量符号 & 再乘功率缩放 \\
\bottomrule
\end{longtable}

\subsection{本模块不负责的内容}

\begin{longtable}{L{5.5cm}L{8.5cm}}
\caption{本模块不负责的内容\label{tab:656-out}} \\
\toprule
\textbf{不负责} & \textbf{负责节} \\
\midrule
CRC / 编码 / 加扰 & 6.5.1、6.2.6、6.5.7 \\
参考信号 QPSK & 6.5.3 \\
RE 映射、功控、MCS 选择 & 6.2.3、6.5.4、6.2.5 \\
\bottomrule
\end{longtable}

\section{在 SLB 通信流程中的角色}

发送链中位于加扰之后、资源映射之前。接收端软解调/LLR 须使用\textbf{同一} Gray 嵌套结构。PHY 不得自主升阶，只执行配置的调制阶。

%----------------------------------------------------------------------
\section{关键参数与强制要求}
%----------------------------------------------------------------------

\textbf{通用约定}：比特按自然顺序消耗；$d(i)$ 对应 $b(mi)\ldots b(mi+m-1)$；I 取偶数下标位、Q 取奇数下标位；公式整体带负号。

\subsection{6.5.6.1 QPSK}

\begin{equation}
d(i)=-\frac{1}{\sqrt{2}}\bigl[(2b(2i)-1)+j(2b(2i+1)-1)\bigr]
\end{equation}

\begin{longtable}{ccccL{5cm}}
\caption{QPSK 真值表（含前置负号）\label{tab:656-qpsk}} \\
\toprule
\textbf{$b(2i)$} & \textbf{$b(2i+1)$} & \textbf{内层} & \textbf{$d(i)$} & \textbf{浮点约值} \\
\midrule
0 & 0 & $-1-j$ & $(+1+j)/\sqrt{2}$ & $(+0.707,+0.707)$ \\
0 & 1 & $-1+j$ & $(+1-j)/\sqrt{2}$ & $(+0.707,-0.707)$ \\
1 & 0 & $+1-j$ & $(-1+j)/\sqrt{2}$ & $(-0.707,+0.707)$ \\
1 & 1 & $+1+j$ & $(-1-j)/\sqrt{2}$ & $(-0.707,-0.707)$ \\
\bottomrule
\end{longtable}

\subsection{16QAM 及以上}

16QAM（$m{=}4$，$1/\sqrt{10}$）示例：
\begin{equation}
\begin{aligned}
d(i)=-\frac{1}{\sqrt{10}}\Big\{
&(2b(4i)-1)\bigl[2-(1-2b(4i+2))\bigr]\\
&+j(2b(4i+1)-1)\bigl[2-(1-2b(4i+3))\bigr]\Big\}
\end{aligned}
\end{equation}

64/256/1024/4096QAM 在 I/Q 上继续嵌套 $4-(1-2b)[\cdots]$、$8-(1-2b)[\cdots]$ 等层，归一化分别为 $1/\sqrt{42}$、$1/\sqrt{170}$、$1/\sqrt{682}$、$1/\sqrt{2730}$。统一递推：
\begin{align}
A_0 &= 2-(1-2b)\in\{1,3\}\\
A_k &= 2^{k+1}-(1-2b_k)A_{k-1}
\end{align}
再乘符号位 $(2b_{\mathrm{sign}}-1)$，整体乘 $-1/\sqrt{N_{\mathrm{avg}}}$。

\begin{longtable}{L{3cm}L{2.5cm}L{8.5cm}}
\caption{符号内比特角色\label{tab:656-bits}} \\
\toprule
\textbf{角色} & \textbf{下标} & \textbf{说明} \\
\midrule
I 符号位 & $b(mi)$ & 首比特 \\
Q 符号位 & $b(mi+1)$ & \\
I 幅度（内$\rightarrow$外） & $b(mi+2),b(mi+4),\ldots$ & \\
Q 幅度（内$\rightarrow$外） & $b(mi+3),b(mi+5),\ldots$ & \\
\bottomrule
\end{longtable}

\textbf{与 6.5.3 对照}：6.5.3 参考 QPSK 为 $+1/\sqrt{2}$（无整体负号）；6.5.6.1 为 $-1/\sqrt{2}$。金标：$b{=}(1,0)$ $\Rightarrow$ $d\approx(-0.707,+0.707)$。

\begin{longtable}{L{4.2cm}L{3cm}L{6.5cm}}
\caption{链路调制选用\label{tab:656-link}} \\
\toprule
\textbf{链路} & \textbf{标准} & \textbf{调制} \\
\midrule
同步/广播/接入/GCI & 6.2.4 & QPSK \\
TCI（含二阶段 ACK） & 6.2.4.10.5 & QPSK 或 16QAM \\
第一/二类数据 & 6.2.5 & QPSK\textasciitilde 4096QAM（MCS 表） \\
\bottomrule
\end{longtable}

TCI 二阶段 HARQ-ACK：G 链路数据为 QPSK/16QAM 时用 QPSK，否则用 16QAM。

%----------------------------------------------------------------------
\section{本节核心详解}
%----------------------------------------------------------------------

\textbf{推荐实现}：离线用标准公式生成 $2^m$ 点 LUT（\texttt{generate\_mod\_lut.py}），运行时按比特打包索引查表；MCS 切换仅换表指针（\texttt{slb\_mod\_select}）。禁止手改生成后的 LUT 源文件。

索引构造（先消耗比特为 MSB）：$\mathrm{idx}=\sum_{j=0}^{m-1} b(mi+j)\,2^{m-1-j}$。

%----------------------------------------------------------------------
\section{数据类型、长度与维度}
%----------------------------------------------------------------------

\begin{longtable}{L{3.5cm}L{2.5cm}L{8cm}}
\caption{调制接口维度\label{tab:656-dim}} \\
\toprule
\textbf{变量} & \textbf{方向} & \textbf{说明} \\
\midrule
\texttt{bits[]} & 输入 & $N_{\mathrm{bit}}$ 个 $\{0,1\}$ \\
\texttt{mod\_type} & 输入 & 六种枚举之一 \\
\texttt{symbols[]} & 输出 & $N_{\mathrm{sym}}=N_{\mathrm{bit}}/m$ 复数 \\
\texttt{m} & 派生 & $\{2,4,6,8,10,12\}$ \\
\bottomrule
\end{longtable}

%----------------------------------------------------------------------
\section{时序关系与触发条件}
%----------------------------------------------------------------------

\begin{longtable}{L{3.5cm}L{5.5cm}L{5.5cm}}
\caption{调用触发条件\label{tab:656-trig}} \\
\toprule
\textbf{能力} & \textbf{进入条件} & \textbf{失败/回退} \\
\midrule
调制映射 & 比特缓冲就绪；数据面已加扰；$N_{\mathrm{bit}}\bmod m{=}0$ & 非整除/非法比特 $\rightarrow$ \texttt{ERR\_PARAM} \\
MCS 换表 & 新 \texttt{mod\_type} 合法 & 非法枚举 $\rightarrow$ 错误，不静默沿用旧表 \\
\bottomrule
\end{longtable}

%----------------------------------------------------------------------
\section{处理步骤}
%----------------------------------------------------------------------

\subsection{端到端：加扰后比特 → 符号 → 映射}

\textbf{Step 1（校验）}
\begin{itemize}[nosep]
  \item \textbf{进入}：\texttt{mod\_type}、$N_{\mathrm{bit}}$、缓冲指针有效；
  \item \textbf{动作}：查 $m$；若 $N_{\mathrm{bit}}\bmod m\neq 0$ 或比特$\notin\{0,1\}$ 则返回错误；
  \item \textbf{失败}：\texttt{SLB\_PHY\_ERR\_PARAM}。
\end{itemize}

\textbf{Step 2（映射）}
\begin{itemize}[nosep]
  \item \textbf{进入}：校验通过；
  \item \textbf{动作}：对 $i{=}0\ldots N_{\mathrm{sym}}{-}1$，取 $m$ 比特 $\rightarrow$ LUT/公式 $\rightarrow$ $d(i)$（含负号与归一化）；
  \item \textbf{失败}：输出容量不足 $\rightarrow$ \texttt{ERR\_BUF}。
\end{itemize}

\textbf{Step 3（交出）}
\begin{itemize}[nosep]
  \item \textbf{进入}：$d[]$ 完成；
  \item \textbf{动作}：交给 6.2.3 映射或 6.5.8；本模块结束；
  \item \textbf{失败}：不得在此回改加扰/MCS。
\end{itemize}

%----------------------------------------------------------------------
\section{协议中允许本模块改变的参数}
%----------------------------------------------------------------------

\begin{longtable}{L{4cm}L{4cm}L{6cm}}
\caption{允许配置的参数\label{tab:656-tunable}} \\
\toprule
\textbf{参数} & \textbf{来源} & \textbf{影响} \\
\midrule
\texttt{mod\_type} & MCS/固定配置 & 选用 6.5.6.x \\
比特序 & 上游编码/加扰 & 必须收发一致 \\
\bottomrule
\end{longtable}

%----------------------------------------------------------------------
\section{约束与实现要点}
%----------------------------------------------------------------------

\begin{enumerate}[nosep]
  \item \textbf{【比特整除】}：$N_{\mathrm{bit}}$ 必须被 $m$ 整除；
  \item \textbf{【前置负号】}：省略负号或套用 6.5.3 符号约定将导致 EVM 失败；
  \item \textbf{【数据先加扰】}：数据面输入必须是 6.5.7 输出，禁止对原始编码比特直接调制；
  \item \textbf{【禁混 RS API】}：参考信号不得调用本节接口；
  \item \textbf{【平均能量 1】}：归一化系数不得改写；
  \item \textbf{【LLR 对称】}：接收软解调须镜像同一 Gray 嵌套；
  \item \textbf{【无自主升阶】}：PHY 只执行配置阶，不根据信道自行改 MCS；
  \item \textbf{【无状态机】}：映射为纯函数（当前比特片 + 调制阶 $\rightarrow$ 符号）；MCS 切换仅换 LUT 指针，不维护 IDLE/BUSY 等 FSM。
\end{enumerate}

%----------------------------------------------------------------------
\section{与标准其他章节的衔接}
%----------------------------------------------------------------------

\begin{itemize}[nosep]
  \item \textbf{6.2.4/6.2.5}：控制固定阶与数据 MCS 表；
  \item \textbf{6.5.7}：数据加扰在调制前；
  \item \textbf{6.5.3}：参考 QPSK 独立实现；
  \item \textbf{6.2.3/6.5.4/6.5.8}：映射、功控、深覆盖 DFT 预编码。
\end{itemize}

\vfill
\begin{center}
\small\textcolor{gray}{精确参数与字段定义以 T/XS 10001-2025 正式标准为准；\\
本文档提供 6.5.6 调制方式工程解读，供 SLB 实现与联调参考。}
\end{center}

\end{document}
